% Partie : théorie de la démonstration
% Chapitre : calcul des séquents

\documentclass[../../../include/open-logic-chapter]{subfiles}

\begin{document}

\olchapter{pt}{seq}{Le calcul des séquents}

\olimport{introduction}
\olimport{rules-proofs}
\olimport{proof-examples}
\olimport{interpretation-rules}
\olimport{quantifiers}
\olimport{admissible-derivable}
\olimport{invertibility}
\olimport{translations}
\begingroup
\olfileid[fr]{pt}{seq}{variants}
\olsection{Tableaux de calculs apparentés}
Les deux premières tables autorisent au plus une formule dans le conséquent. Les deux suivantes en admettent plusieurs.
\olimport{rules-G1i}
\olimport{rules-G3i}
\olimport{rules-LK}
\olimport{rules-mG3i}
\endgroup
\clearpage
\OLEndChapterHook

\end{document}
